% Nouveau EGA V (cote 25): archive cote 25, site page 120 (PDF page 121).
% typescript. Draft compared: none.
% AI-transcribed by Claude Opus 5.5 (claude-opus-5-5, adaptive thinking, effort high), Message Batches API, batch ega25b (msgbatch_01A7n5YuXubUPNK1zsepi8RG), 2026-11-13; not yet checked by eye.
\scanpage{25/120}
\mspage{\unsure{27}}

Z$_s$ est de codimension $\geq 2$ dans X$_s$ ; \xout{\illegible{2}\unsure{(comme dans les cas envisa}}
\xout{familles} \inserted{syst lin} de diviseurs par.par P qui so\inserted{\unsure{ie}}nt
\uline{sans composante fixe}, \xout{tandis que $\emptyset$ est l'ensemble de} D'autre part, si
S=Spec(k), et si X est normal, alors \inserted{pour} \uline{toute} application rationelle
f:X$\to$P, l'ensemble de définition U(f) vérifie codim ((X-U(f),X)$\geq$2
(II 7), donc dans ce cas $\emptyset$ est formé de l'ensemble de \uline{toutes} les
applications rationelles de X dans P.

La démonstration de a) a déjà été donnée. Pour prouver b), notons
que la formation de $i_*(\underline{L}_U)$ commute avec toute extension plate de la
\xout{pour le\unsure{x} il faut} \inserted[typed]{la cond. envisagée est inv. par}
base S' \quad S (réf)\inserted{\unsure{du moins si $U\to X$ est quasi-compact, cas auquel on se ramène sans difficultés}}, donc \xout{il suffit de trouver un morphisme} fid plat qu cpc
\xout{S' $\quad$ S tel que la conclusion soit vraie après} changement de base. On
prendra S' = $\hand{\check{}}\mathrm{P}$ , et on note que l'hypothèse $i'_*(\underline{L}'_U)$ inversible ne
change pas en remplaçant $\underline{L}'_U$ par $\underline{L}'_U\otimes_{S'}M'$, où M' est un Module inversible
sur S' (car $i'_*(\underline{L}'_U\otimes_{S'}M')\simeq i'_*(\underline{L}'_U)\otimes_{S'}M'$ ). On \unsure{prèndra|prendra} M'= $\underline{O}_{\check{P}}(1)$, de
sorte que la condition envisagée signifie aussi que $i'_*(\underline{N}')$ est un Module
inversible, où $\underline{N}'= (f_U\times_S id_P)^*(\underline{O}_{P\times\check{P}}(1,1))$ . Or $\underline{O}$ (1,1) est précisément
le Module inversible défini par le diviseur canonique $\underline{H}$ sur $P\times_S\check{P}$ , de sor%
te que $\underline{N}$ n'est autre que le Module inversible défini par \qquad D' =
$(f_U\times id_P)$ $(\underline{H})$. \deleted{\xout{\illegible{1}}} Lorsque f provient de D$\in\emptyset$ , \xout{ce dernier}
D' n'est autre que D$|$U$\times_S\check{P}$ , \xout{et il résulte du lemme} donc $\underline{N}' = \underline{N}|U\times_S\check{P}$ ,
où $\underline{N}$ est le Module inversible sur $X\times_S\hand{\check{}}P$ défini par D, et il résulte
du lemme \inserted{3.2} plus haut \xout{\illegible{1}} appliqué à $P\times_S\hand{\check{}}P\to\hand{\check{}}P$ qu'on aura
$i'_*(\underline{N}')\simeq\underline{N}$ , donc $i'_*(\underline{N}')$ est inversible. Inversement, si cette condi-%
tion est remplie prouvons que f provient d'un\xout{e} D$\in\emptyset$ , ou ce qui revient
ma\xout{n}nifestement au même, que le diviseur D' se prolonge en un diviseur \deleted{de}
\deleted{Cartier} relatif par r apport \xout{de} à $\hand{\check{}}P$ sur X$\times_S$P . Il \deleted{suffit} revient au même
de dire qu'il se prolonge en un diviseur \deleted{de Cartier} D sur X$\times_S$P , car \hand{D}.

\margin{\unsure{Rq compliquée} \unsure{Il faut expliquer} \illegible{2} \unsure{$f\times\mathrm{id}$} \unsure{si $p^{-1}(H)$} \unsure{est défini} \illegible{1} \unsure{diviseur rel.} \unsure{à la condition} \illegible{3} \unsure{diviseur relatif} \unsure{$X\times P$} \unsure{/P} \illegible{2}}

\mspage{120}
